% GENERATED FILE. Edit canonical lessons, then run npm run generate:books.
% canonical-source-hash: fnv1a64:f967571b2b239492
% canonical-lessons: ZH-C316-dengji, ZH-C316-shuxi, ZH-C316-dangdi, ZH-C316-rensheng, ZH-C316-ziran

\chapter{A New Place}
\label{ch:zh-a-new-place}
\hlchaptermodality{316}

\begin{chapteropening}
\textbf{By the end of this chapter:} \emph{I can talk about registering in a new town, getting to know it, and life there.}

Complete the last lesson of chapter 316: I can talk about registering in a new town, getting to know it, and life there.
\end{chapteropening}

\section[dēngjì]{dēngjì — to register officially}

\label{lesson:ZH-C316-dengji}



\begin{quote}
Before the new one: say the Mandarin for to my surprise, then the Mandarin for the same as before.
\end{quote}

\subsection*{You'll want to know: dēngjì}
\emph{dēngjì} — \textquotedblleft{}to register officially\textquotedblright{}.

\begin{quote}
\emph{wǒmen qù pàichūsuǒ dēngjì} — We went to the police station to register.
\end{quote}

Its characters come later; for now it is a word for the ear.

\begin{cousinweb}
\emph{dēngjì} is built from \emph{dēng} 'to go up' and \zh{记} 'to note down', both inherited from Middle Chinese, \emph{tong} and \emph{kiH} (Guangyun, in Baxter's transcription), which glosses the first 'to rise' and 'to go forward': to enter in the record.
\end{cousinweb}

\begin{grammarlens}[title={What you've built}]
One more word for telling what happened.
\end{grammarlens}

\subsection*{Your turn}
\begin{itemize}
\raggedright
  \item \emph{Say it:} \emph{dēngjì}
  \item \emph{Say it:} \emph{dēngjì}, once more
  \item \emph{Say it:} \emph{wǒmen qù pàichūsuǒ dēngjì} — We went to the police station to register.
  \item \emph{Recall:} say the Mandarin for to my surprise, then the Mandarin for the same as before, then say \emph{dēngjì} again
\end{itemize}

\begin{tcolorbox}[breakable,colback=teal!4,colframe=teal!35!black,title={Before you move on}]
What does dēngjì mean? (\textquotedblleft{}To register officially\textquotedblright{}.) Say it once more.
\end{tcolorbox}
\section[shúxī]{shúxī — to know well; to be familiar with}

\label{lesson:ZH-C316-shuxi}



\begin{quote}
Before the new one: say the Mandarin for the same as before, then the Mandarin for to register officially.
\end{quote}

\subsection*{You'll want to know: shúxī}
\emph{shúxī} — \textquotedblleft{}to know well; to be familiar with\textquotedblright{}.

\begin{quote}
\emph{wǒ duì zhèr hái bù shúxī} — I don't know this place well yet.
\end{quote}

Its characters come later; for now it is a word for the ear.

\begin{cousinweb}
\emph{shúxī} is built from \emph{shú} 'ripe' and \emph{xī} 'in full', both inherited from Middle Chinese, \emph{dzyuwk} and \emph{sit} (Guangyun, in Baxter's transcription), which glosses the first 'done' and the second 'in every detail'.
\end{cousinweb}

\begin{grammarlens}[title={What you've built}]
One more word for telling what happened.
\end{grammarlens}

\subsection*{Your turn}
\begin{itemize}
\raggedright
  \item \emph{Say it:} \emph{shúxī}
  \item \emph{Say it:} \emph{shúxī}, once more
  \item \emph{Say it:} \emph{wǒ duì zhèr hái bù shúxī} — I don't know this place well yet.
  \item \emph{Recall:} say the Mandarin for the same as before, then the Mandarin for to register officially, then say \emph{shúxī} again
\end{itemize}

\begin{tcolorbox}[breakable,colback=teal!4,colframe=teal!35!black,title={Before you move on}]
What does shúxī mean? (\textquotedblleft{}To know well; to be familiar with\textquotedblright{}.) Say it once more.
\end{tcolorbox}
\section[\texorpdfstring{\zh{当地} (dāngdì)}{当地 (dāngdì)}]{\zh{当地} (dāngdì) — of that place; in the place in question}

\label{lesson:ZH-C316-dangdi}



\begin{quote}
Before the new one: say the Mandarin for to register officially, then the Mandarin for to be familiar with.
\end{quote}

\subsection*{You'll want to know: \zh{当地}}
\textbf{\zh{当地}} — \emph{dāngdì} — \textquotedblleft{}of that place; in the place in question\textquotedblright{}.

\begin{quote}
\textbf{\zh{我到了当地才找到房子}} \emph{wǒ dào le dāngdì cái zhǎodào fángzi} — I only found somewhere to live once I got there.
\end{quote}

\begin{cousinweb}
\textbf{\zh{当地}} (\emph{dāngdì}) is built from \zh{当} 'at' and \zh{地} 'a place', both inherited from Middle Chinese, \emph{tang} and \emph{dijH} (Guangyun, in Baxter's transcription), which glosses the first 'to face' and 'to meet': the very place.
\end{cousinweb}

\begin{grammarlens}[title={What you've built}]
One more word for telling what happened.
\end{grammarlens}

\subsection*{Your turn}
\begin{itemize}
\raggedright
  \item \emph{Say it:} \emph{dāngdì}
  \item \emph{Say it:} \emph{dāngdì}, once more
  \item \emph{Say it:} \emph{wǒ dào le dāngdì cái zhǎodào fángzi} — I only found somewhere to live once I got there.
  \item \emph{Recall:} say the Mandarin for to register officially, then the Mandarin for to be familiar with, then say \emph{dāngdì} again
\end{itemize}

\begin{tcolorbox}[breakable,colback=teal!4,colframe=teal!35!black,title={Before you move on}]
What does \zh{当地} mean? (\textquotedblleft{}Of that place; in the place in question\textquotedblright{}.) Say it once more.
\end{tcolorbox}
\section[\texorpdfstring{\zh{人生} (rénshēng)}{人生 (rénshēng)}]{\zh{人生} (rénshēng) — a person's life; life from birth to death}

\label{lesson:ZH-C316-rensheng}



\begin{quote}
Before the new one: say the Mandarin for to be familiar with, then the Mandarin for of that place.
\end{quote}

\subsection*{You'll want to know: \zh{人生}}
\textbf{\zh{人生}} — \emph{rénshēng} — \textquotedblleft{}a person's life; life from birth to death\textquotedblright{}.

\begin{quote}
\emph{rénshēng hěn cháng, nǐ kěyǐ chóngxīn kāishǐ} — Life is long; you can start over.
\end{quote}

\begin{cousinweb}
\textbf{\zh{人生}} (\emph{rénshēng}) is built from \zh{人} 'a person' and \zh{生} 'to live', both inherited from Middle Chinese, \emph{nyin} and \emph{srjaeng} (Guangyun, in Baxter's transcription), which glosses the second 'to grow'.
\end{cousinweb}

\begin{grammarlens}[title={What you've built}]
One more word for telling what happened.
\end{grammarlens}

\subsection*{Your turn}
\begin{itemize}
\raggedright
  \item \emph{Say it:} \emph{rénshēng}
  \item \emph{Say it:} \emph{rénshēng}, once more
  \item \emph{Say it:} \emph{rénshēng hěn cháng, nǐ kěyǐ chóngxīn kāishǐ} — Life is long; you can start over.
  \item \emph{Recall:} say the Mandarin for to be familiar with, then the Mandarin for of that place, then say \emph{rénshēng} again
\end{itemize}

\begin{tcolorbox}[breakable,colback=teal!4,colframe=teal!35!black,title={Before you move on}]
What does \zh{人生} mean? (\textquotedblleft{}A person's life; life from birth to death\textquotedblright{}.) Say it once more.
\end{tcolorbox}
\section[\texorpdfstring{\zh{自然} (zìrán)}{自然 (zìrán)}]{\zh{自然} (zìrán) — naturally; nature}

\label{lesson:ZH-C316-ziran}



\begin{quote}
Before the new one: say the Mandarin for of that place, then the Mandarin for life from birth to death.
\end{quote}

\subsection*{You'll want to know: \zh{自然}}
\textbf{\zh{自然}} — \emph{zìrán} — \textquotedblleft{}naturally; nature\textquotedblright{}.

\begin{quote}
\textbf{\zh{又下雨了我们自然不出门}} \emph{yòu xiàyǔ le, wǒmen zìrán bù chūmén} — It's raining again, so naturally we're staying in.
\end{quote}

\begin{cousinweb}
\textbf{\zh{自然}} (\emph{zìrán}) is built from \zh{自} and \zh{然} 'so', both inherited from Middle Chinese, \emph{dzijH} and \emph{nyen} (Guangyun, in Baxter's transcription), which glosses the first 'from' and 'by way of'; its sense 'by itself' is not given there. So of itself.
\end{cousinweb}

\begin{grammarlens}[title={What you've built}]
One more word for telling what happened.
\end{grammarlens}

\subsection*{Your turn}
\begin{itemize}
\raggedright
  \item \emph{Say it:} \emph{zìrán}
  \item \emph{Say it:} \emph{zìrán}, once more
  \item \emph{Say it:} \emph{yòu xiàyǔ le, wǒmen zìrán bù chūmén} — It's raining again, so naturally we're staying in.
  \item \emph{Recall:} say the Mandarin for of that place, then the Mandarin for life from birth to death, then say \emph{zìrán} again
\end{itemize}

\begin{tcolorbox}[breakable,colback=teal!4,colframe=teal!35!black,title={Before you move on}]
What does \zh{自然} mean? (\textquotedblleft{}Naturally; nature\textquotedblright{}.) Say it once more.
\end{tcolorbox}
